\begin{proposition}[Analytic base change in degree two]
\hypertarget{tag-03II}{}\label{wbII:prop:analytic-base-change}
\StatementRecord{wbII:prop:analytic-base-change}{proved}{03II}{\cite[pp.~120--121]{BS76push}; \cite[Lemma~1.5]{CDP98}}

Let $f:X\to B$ be proper between complex spaces, let
$\mathcal F$ be coherent and $f$-flat, and fix $c\in B$.  If
\[
H^3(X_c,\mathcal F|_{X_c})=0,
\tag{W.13}\hypertarget{tag-03ED}{}\label{wbII:eq:W:5.16}
\]
then the canonical fibre map
\[
\beta_c^2:
(R^2f_*\mathcal F)_c\otimes_{\mathcal O_{B,c}}\mathbb C(c)
\longrightarrow H^2(X_c,\mathcal F|_{X_c})
\tag{W.14}\hypertarget{tag-03EE}{}\label{wbII:eq:W:5.17}
\]
is an isomorphism.
\end{proposition}



\begin{proof}
Apply B\u{a}nic\u{a}--St\u{a}n\u{a}\c{s}il\u{a},
Chapter~III, \S3, Corollary~3.5, with $q=3$
\cite[pp.~120--121]{BS76push}.  Its hypotheses are precisely
properness, coherence, and flatness.  Condition \eqref{wbII:eq:W:5.16} first forces the
degree-three direct image to vanish near $c$.  The corollary then applies the
equivalent base-change conditions of its preceding Theorem~3.4 in degree
$q-1=2$ and states that the canonical map \eqref{wbII:eq:W:5.17} is bijective.  The top-degree application also appears in the proof of
\cite[Lemma~1.5(c), p.~80]{CDP98}.

\end{proof}